% Part: incompleteness
% Chapter: theories-computability
% Section: first-incompleteness

\documentclass[../../../include/open-logic-section]{subfiles}

\begin{document}

\olfileid{inc}{tcp}{inc}

\olsection{$\Th{Q}$కు సంపూర్ణమైన, అవైరుధ్యమైన,
  \tetoken{స్వీకృతీకరించదగిన}{!!^{axiomatizable}} విస్తరణలు లేవు}

\begin{thm}
\ollabel{thm:first-incompleteness}
$\Th{Q}$కు సంపూర్ణమైన, అవైరుధ్యమైన, \tetoken{స్వీకృతీకరించదగిన}{!!{axiomatizable}} విస్తరణ ఏదీ లేదు.
\end{thm}

\begin{proof}
$\Th{Q}$కు అవైరుధ్యమైన, నిర్ణయించదగిన విస్తరణ ఏదీ లేదని మనకు ఇప్పటికే
తెలుసు. కానీ $\Th{T}$ సంపూర్ణమూ \tetoken{స్వీకృతీకరించదగినదీ}{!!{axiomatized}} అయితే, అది
నిర్ణయించదగినది.\sourcecorrection{OLTETCP-001}{ఆంగ్ల మూల నిరూపణ ``axiomatized'' అని మాత్రమే చెప్పింది; దాని ముందు ఉన్న ఉపసిద్ధాంతమూ ప్రస్తుత సిద్ధాంతపు పరికల్పనమూ గణనీయ స్వీకృతాల సమితి గల ``axiomatizable'' సిద్ధాంతాన్నే కోరుతాయి. ఆ అవసరమైన అర్థాన్ని అనువాదంలో నిలిపాం.}
\end{proof}

\begin{explain}
ఈ సిద్ధాంతం గ్యోడెల్ 1931లో ఇచ్చిన మొదటి అసంపూర్ణతా సిద్ధాంతపు అసలు
సూత్రీకరణకు చాలా దూరంలో లేదు. మరింత ఆధునిక పదజాలాన్ని పక్కన పెడితే, ప్రధాన
తేడాలు ఇవి: గ్యోడెల్ ``అవైరుధ్యమైన'' స్థానంలో ``$\omega$-అవైరుధ్యమైన'' అని
వాడాడు; అలాగే గణనీయతకు అధికారిక భావన అప్పటికి సిద్ధంగా లేకపోవడం వల్ల,
``\tetoken{స్వీకృతీకరించదగిన}{!!{axiomatizable}}'' అని పూర్తి సాధారణతతో చెప్పలేకపోయాడు. (తదుపరి
దశాబ్దంలో గ్యోడెల్ సహా పలువురు గణనీయతకు అధికారిక నమూనాలను అభివృద్ధి చేశారు;
గ్యోడెల్ ఘటనకు లోబడే సిద్ధాంతాల రకాలను లక్షణీకరించగలగడమే ఆ పనికి చాలావరకు
ప్రేరణ.)

సంపూర్ణత, అవైరుధ్యం, \tetoken{స్వీకృతీకరించదగినత}{!!{axiomatizability}}---ఈ మూడూ ఒకేసారి ఉండలేవని
సిద్ధాంతం చెబుతుంది. అయితే వీటిలో ఏదో ఒకదాన్ని వదిలేస్తే మిగిలిన రెండూ
ఉండవచ్చు: $\Th{Q}$ అవైరుధ్యమైనది, గణనీయంగా స్వీకృతీకరించబడినది, కాని
సంపూర్ణం కాదు; వైరుధ్య సిద్ధాంతం సంపూర్ణం, గణనీయంగా స్వీకృతీకరించబడినది
(ఉదాహరణకు $\{ \eq/[0][0] \}$ చేత), కాని అవైరుధ్యమైనది కాదు; అంకగణితపు
సత్య వాక్యాల సమితి సంపూర్ణమూ అవైరుధ్యమూ, కాని అది గణనీయంగా
స్వీకృతీకరించబడినది కాదు.
\end{explain}
\end{document}
